\documentclass[letterpaper]{article} % DO NOT CHANGE THIS
\usepackage[preprint]{aaai2027}  % DO NOT CHANGE THIS
% The serif, sans-serif, and monospaced fonts are loaded automatically by
% aaai2027.sty (newtxtext, helvet, courier). DO NOT add \usepackage{times},
% \usepackage{helvet}, \usepackage{courier}, or any other font package.
\usepackage[hyphens]{url}  % DO NOT CHANGE THIS
\usepackage{graphicx} % DO NOT CHANGE THIS
\usepackage{enumitem}
\urlstyle{rm} % DO NOT CHANGE THIS
\def\UrlFont{\rm}  % DO NOT CHANGE THIS
\usepackage{natbib}  % DO NOT CHANGE THIS AND DO NOT ADD ANY OPTIONS TO IT
\usepackage{caption} % DO NOT CHANGE THIS AND DO NOT ADD ANY OPTIONS TO IT
\frenchspacing  % DO NOT CHANGE THIS
%
% These are recommended to typeset algorithms but not required. See the subsubsection on algorithms. Remove them if you don't have algorithms in your paper.
\usepackage{algorithm}
\usepackage{algorithmic}
\usepackage{amsmath}
\usepackage{amssymb}
%
% These are recommended to typeset listings but not required. See the subsubsection on listing. Remove this block if you don't have listings in your paper.
\usepackage{newfloat}
\usepackage{listings}
\DeclareCaptionStyle{ruled}{labelfont=normalfont,labelsep=colon,strut=off} % DO NOT CHANGE THIS
\lstset{%
	basicstyle={\footnotesize\ttfamily},% footnotesize acceptable for monospace
	numbers=left,numberstyle=\footnotesize,xleftmargin=2em,% show line numbers, remove this entire line if you don't want the numbers.
	aboveskip=0pt,belowskip=0pt,%
	showstringspaces=false,tabsize=2,breaklines=true}
\floatstyle{ruled}
\newfloat{listing}{tb}{lst}{}
\floatname{listing}{Listing}

%
% Recommended for better-looking tables
\usepackage{booktabs}

%
% Keep the \pdfinfo as shown here. There's no need
% for you to add the /Title and /Author tags.
\pdfinfo{
/TemplateVersion (2027.1)
}

% DISALLOWED PACKAGES
% \usepackage{authblk} -- This package is specifically forbidden
% \usepackage{balance} -- This package is specifically forbidden
% \usepackage{CJK} -- This package is specifically forbidden
% \usepackage{float} -- This package is specifically forbidden
% \usepackage{flushend} -- This package is specifically forbidden
% \usepackage{fullpage} -- This package is specifically forbidden
% \usepackage{geometry} -- This package is specifically forbidden
% \usepackage{hyperref} -- This package is specifically forbidden
% \usepackage{navigator} -- This package is specifically forbidden
% (or any other package that embeds links such as navigator or hyperref)
% \usepackage{indentfirst} -- This package is specifically forbidden
% \usepackage{layout} -- This package is specifically forbidden
% \usepackage{multicol} -- This package is specifically forbidden
% \usepackage{nameref} -- This package is specifically forbidden
% \usepackage{savetrees} -- This package is specifically forbidden
% \usepackage{setspace} -- This package is specifically forbidden
% \usepackage{stfloats} -- This package is specifically forbidden
% \usepackage{tabu} -- This package is specifically forbidden
% \usepackage{titlesec} -- This package is specifically forbidden
% \usepackage{tocbibind} -- This package is specifically forbidden
% \usepackage{ulem} -- This package is specifically forbidden
% \usepackage{wrapfig} -- This package is specifically forbidden
% DISALLOWED COMMANDS
% \nocopyright -- Your paper will not be published if you use this command
% \addtolength -- This command may not be used
% \balance -- This command may not be used
% \baselinestretch -- Your paper will not be published if you use this command
% \clearpage -- No page breaks of any kind may be used for the final version of your paper
% \columnsep -- This command may not be used
% \newpage -- No page breaks of any kind may be used for the final version of your paper
% \pagebreak -- No page breaks of any kind may be used for the final version of your paper
% \pagestyle -- This command may not be used
% \tiny -- This is not an acceptable font size.
% \vspace{- -- No negative value may be used in proximity of a caption, figure, table, section, subsection, subsubsection, or reference
% \vskip{- -- No negative value may be used to alter spacing above or below a caption, figure, table, section, subsection, subsubsection, or reference

\setcounter{secnumdepth}{0} %May be changed to 1 or 2 if section numbers are desired.

% The file aaai2027.sty is the style file for AAAI Press
% proceedings, working notes, and technical reports.
%

% Title

% Your title must be in mixed case, not sentence case.
% That means all verbs (including short verbs like be, is, using,and go),
% nouns, adverbs, adjectives should be capitalized, including both words in hyphenated terms, while
% articles, conjunctions, and prepositions are lower case unless they
% directly follow a colon or long dash
\title{BRA-Audit: Budgeted Runtime Auditing for LLM Multi-Agent Systems via Cumulative-Exposure Audit-Point Placement}
\author{
    %Authors
    % All authors must be in the same font size and format.
    Kaixiang Wang,
    Yidan Lin,
    Jiong Lou,
    Jie Li
}
\affiliations{
    %Afiliations
    
% See more examples next
}

%Example, Single Author, ->> remove \iffalse,\fi and place them surrounding AAAI title to use it
\iffalse
\title{My Publication Title --- Single Author}
\author {
    Author Name
}
\affiliations{
    Affiliation\\
    Affiliation Line 2\\
    name@example.com
}
\fi

\iffalse
%Example, Multiple Authors, ->> remove \iffalse,\fi and place them surrounding AAAI title to use it
\title{My Publicatio}
\author {
    % Authors
    First Author Name\textsuperscript{\rm 1,\rm 2}\equalcontrib,
    Second Author Name\textsuperscript{\rm 2}\equalcontrib,
    Third Author Name\textsuperscript{\rm 1}\corresponding
}
\affiliations {
    % Affiliations
    \textsuperscript{\rm 1}Affiliation 1\\
    \textsuperscript{\rm 2}Affiliation 2\\
    firstAuthor@affiliation1.com, secondAuthor@affilation2.com, thirdAuthor@affiliation1.com
}
\fi


\begin{document}

\maketitle

\begin{abstract}
% Large language model-based multi-agent systems (LLM-MAS) enable complex problem solving through role specialization, communication, and collaborative reasoning. Yet their inter-agent dependencies also make execution fragile: hallucinated or malicious outputs from only a few agents can propagate through the workflow and cause system-level failures. Recent safeguards often introduce dedicated auditor agents, but they face a fundamental efficiency dilemma. Auditing only at the end requires reasoning over long and complex trajectories, which weakens detection and enlarges rollback scope. Auditing every agent at every round enables strong detection performance and precise localization, but repeated auditor calls incur high token cost. 

% Large language model-based multi-agent systems (LLM-MAS) enable complex problem solving through specialized collaboration, but their inter-agent dependencies allow hallucinated or malicious outputs to propagate into system-level failures. Dedicated auditor agents can mitigate these risks, yet existing strategies face an efficiency dilemma: end-only auditing requires processing long trajectories, which weakens audit effectiveness and leads to coarse-grained rollback, while auditing every agent at every round improves detection and localization at substantial token cost. To address this gap, we propose BRA-Audit, a budget-aware runtime auditing framework. BRA models MAS execution as a dynamic dependency graph and formulates audit scheduling as a audit point-placement problem that minimizes cumulative unchecked exposure. It prioritizes influential and long-unaudited regions and supports localized recovery from trusted audit points. Experiments across MAS collaboration environments and complex open-ended tasks show that BRA-Audit restores performance close to the clean setting, achieves results comparable to PeerGuard and Multi-Agent Debate, and reduces token consumption by \(17.2\%\)--\(40.6\%\). Our work is available on \url{https://anonymous.4open.science/r/BRA-7CC6/}.

% Large language model-based multi-agent systems (LLM-MAS) enable complex problem solving through specialized collaboration, but their inter-agent dependencies allow hallucinated or malicious outputs to propagate into system-level failures. Dedicated auditor agents can mitigate these risks, yet existing strategies face an efficiency dilemma: end-only auditing examines the entire MAS execution trajectory together with the final output, resulting in an excessively long audit context that may degrade audit effectiveness and increase the scope and cost of rollback, whereas auditing every agent at each round improves detection and localization at substantial token cost. This raises a key question: how can we preserve guard performance while minimizing token cost? 


LLM-based multi-agent systems (LLM-MAS) solve complex tasks through specialized collaboration, but inter-agent dependencies can propagate hallucinated or malicious outputs into system-level failures. Auditor agents mitigate these risks, yet existing strategies face an efficiency dilemma: end-only auditing reviews long trajectories and final outputs, potentially weakening audit effectiveness and enlarging rollback scope, while auditing every agent each round improves detection and localization at high token cost. How can guard performance be preserved while minimizing token cost?
To address this problem, we propose BRA-Audit, a budget-aware runtime auditing framework that models MAS execution as a dynamic dependency graph and formulates audit scheduling as audit-point placement under a fixed audit-call budget to minimize cumulative unchecked exposure. Its greedy scheduler prioritizes influential and long-unaudited regions, while trusted audit points enable localized recovery. Across structured coordination, complex reasoning, and open-ended tasks, BRA-Audit restores performance close to the clean setting, remains competitive with heavy guard methods and reduces end-to-end token consumption by \(17.2\%\)--\(40.6\%\). 
\end{abstract}



\section{Introduction}
LLM-based multi-agent systems (LLM-MAS) coordinate specialized agents through explicit roles, communication, and shared workflows~\cite{hong2024metagpt,qian2024chatdev,chan2023chateval}. Recent systems demonstrate that MAS can support generalist task execution, diverse coordination protocols, multi-agent reflection, and layered response aggregation~\cite{fourney2024magenticonegeneralistmultiagentsolving,zhu2025multiagentbenchevaluatingcollaborationcompetition,wang2024mixtureofagentsenhanceslargelanguage}. Other studies optimize communication graphs and architecture selection to remove redundant exchanges and adapt computation to task difficulty~\cite{zhang2024cutcrapeconomicalcommunication,li2024survey,guo2024large}.
These advances suggest that LLM-MAS can combine task decomposition, heterogeneous expertise, tool use, and multiple reasoning paths, making them promising for long-horizon and complex open-ended tasks~\cite{tran2025collaboration,yan2026selftalk,li2023camel}.

However, multi-agent coordination also expands the failure surface. Hallucinated or malicious messages can be reused as context, alter downstream decisions, and propagate across agents~\cite{lee2024promptinfection,ju2024flooding,he2025redteaming}. Recent attacks can manipulate exchanged messages, hijack execution control, leak system information, or conceal malicious intent~\cite{triedman2025maliciouscode,wang2025masleak,xie2025mole}. Because agents repeatedly consume one another's outputs, a local error may evolve into a system-level failure, while its origin becomes difficult to identify in long and tightly coupled trajectories~\cite{xia2026agentlocate,zheng2025integrity}.

\begin{figure}[t]
    \centering
    \includegraphics[width=\linewidth]{Figures/pareto.pdf}
    \caption{The performance-cost trade-off of BRA-Audit against other methods under attack in AgentsNet benchmark.}
    \label{fig:pareto}
\end{figure}

% Early safeguards mainly rely on lightweight mechanisms. Prompt-based self-reflection asks an agent to critique its own output, but the reviewer may share the same failure mode as the generator
% \cite{renze2024selfreflection}. G-Safeguard detects anomalous agents from the communication graph, yet its supervised detector relies on labeled attacks and may generalize poorly to unseen threats~\cite{wang2025gsafeguard,miao2025blindguard}. To improve reliability, recent methods introduce dedicated reviewer agents. GuardAgent and TrinityGuard perform explicit runtime checks, while PeerGuard and multi-agent debate use cross-agent reasoning to identify unreliable outputs~\cite{xiang2024guardagent,wang2026trinityguard,
% fan2025peerguard,du2024multiagentdebate}.

% Existing safeguards mainly focus on how to perform auditing, while paying limited attention to when and where audits should be placed. Auditing only at the end requires the reviewer to process a long and complex interaction trajectory. The extended context and dense dependencies increase the reasoning burden, reduce audit accuracy, and make failure localization difficult. Once an error is detected, the system may also need to roll back a large portion of the execution~\cite{liu-etal-2024-lost,yang2026agentauditor}. In contrast, auditing
% every agent at every round enables accurate detection and fine-grained localization, but requires repeated calls to the auditor agent, substantially increasing token cost~\cite{fan2025peerguard,du2024multiagentdebate}. Since most benigninteractions require no intervention, exhaustive auditing also wastes considerable resources. \textbf{This creates a fundamental dilemma: sparse auditing reduces runtime cost but weakens detection and localization, whereas dense auditing provides precise oversight at a high cost.}

% Early safeguards use lightweight mechanisms. Prompt-based
% self-reflection asks an agent to critique its own output, but the
% reviewer may share the generator's failure modes
% \cite{shinn2023reflexion,renze2024selfreflection}. G-Safeguard detects anomalous agents from communication graphs, yet relies on labeled attacks and may generalize poorly to unseen threats~\cite{wang2025gsafeguard,miao2025blindguard}. To improve reliability,
% recent methods introduce dedicated reviewers. GuardAgent and
% TrinityGuard perform runtime checks, while PeerGuard and multi-agent
% debate use cross-agent reasoning to detect unreliable outputs
% \cite{xiang2024guardagent,wang2026trinityguard,
% fan2025peerguard,du2024multiagentdebate}.

% Most safeguards focus on how to audit, with less attention to when and
% where audits should occur. End-only auditing requires processing a long,
% dependency-dense trajectory, which can reduce audit accuracy, hinder
% failure localization, and require large rollback cost
% \cite{liu-etal-2024-lost,yang2026agentauditor} (As shown in Fig.~\ref{fig:audit_dilemma}a). Auditing every agent at every round provides stronger detection performance and finer localization, but
% repeated auditor calls substantially increase cost~\cite{fan2025peerguard,du2024multiagentdebate}. Since most benign
% interactions require no intervention, exhaustive auditing is often
% wasteful (As shown in Fig.~\ref{fig:audit_dilemma}b). This creates a significant dilemma: sparse auditing
% reduces cost but weakens detection performance and localization, whereas dense auditing provides precise oversight at a high cost.


Past safeguards detect unreliable behavior through self-reflection or graph-based anomaly detection. However, self-reflection may reproduce the generator's failure modes, while GNN-based detectors often rely on labeled attacks and generalize poorly~\cite{shinn2023reflexion,renze2024selfreflection,wang2025gsafeguard,miao2025blindguard}. Dedicated auditor agents have therefore become a common alternative, providing stronger and more transferable oversight through runtime inspection and cross-agent reasoning~\cite{xiang2024guardagent,wang2026trinityguard,fan2025peerguard,du2024multiagentdebate}. Nevertheless, existing studies mainly improve \emph{how} auditing is performed, with limited attention to \emph{when} and \emph{where} auditors should be invoked. End-only auditing examines the full execution trajectory and final result after task completion. Its long, dependency-dense context can obscure failure sources, weaken audit effectiveness, delay intervention, and increase rollback costs~\cite{liu-etal-2024-lost,yang2026agentauditor}. In contrast, interaction-level auditing, as exemplified by PeerGuard~\cite{fan2025peerguard}, inspects agents at each round and enables earlier detection and finer localization. However, repeated auditor calls incur substantial token costs and waste computation on benign interactions~\cite{du2024multiagentdebate}. This raises a significant question: \emph{how can we determine when and where to audit while preserving guard performance and minimizing token cost?}

To answer this  question, we propose BRA-Audit, a budget-aware runtime
framework that directs limited audit resources to critical points in LLM-MAS execution. BRA models inter-agent communication and execution dependencies as a dynamic graph, where nodes represent runtime outputs and edges capture information flow. It then formulates audit scheduling as an audit-point placement problem that minimizes cumulative unchecked exposure under a fixed budget. The scheduler combines audit gaps with topological influence, prioritizing regions that have remained unchecked and can strongly affect downstream computation. Each verified audit point becomes a trusted audit point in this round; when unsafe behavior is detected, BRA contains the affected region and rolls back only the necessary subtrajectory. By integrating topology-aware scheduling, bounded auditing, and localized recovery, BRA maintains effective runtime oversight while reducing redundant checks and re-execution. It therefore achieves an empirical trade-off between audit performance and token consumption.

Further experimental results show that: across MAS collaboration environments and complex open-ended tasks, BRA-Audit restores task performance close to clean baselines while remaining competitive with stronger auditing baselines. It reduces
token consumption by \(17.2\%\)--\(40.6\%\) compared with heavy guard like PeerGuard and Multi-Agent Debate, demonstrating a favorable trade-off between defense effectiveness and token cost (as shown in Fig.~\ref{fig:pareto}).

Our contributions are as follows:
\begin{itemize}[leftmargin=2em]
    \item To determine when and where to audit an LLM-MAS while minimizing cost without compromising audit effectiveness, we formulate budget-constrained auditing as an audit-point placement problem over dynamic execution graphs.
    
    
    
    % We formulate budget-constrained auditing for LLM-MAS as an audit-point placement problem over dynamic execution graphs, capturing the trade-off among detection effectiveness, audit cost, and rollback scope.

    \item We propose BRA-Audit, an online budgeted audit-point placement framework in LLM-MAS that sequentially selects audit points by marginal cumulative-exposure reduction. It combines audit gaps and downstream dependency reachability to weight unchecked exposure, while verified audit points bound audit contexts and localize rollback.
    
    % We propose BRA-Audit, a topology-aware auditing framework that prioritizes influential and long-unaudited regions, limits cumulative unchecked exposure, and supports localized recovery from trusted audit points.

    \item We conduct extensive experiments in MAS collaboration environments and complex open-ended tasks. BRA-Audit achieves competitive recovery performance while reducing token consumption by \(17.2\%\)--\(40.6\%\) compared with strong auditing baselines, demonstrating a favorable empirical Pareto trade-off.
\end{itemize}

\section{Related Work}


% \paragraph{LLM Multi-Agent Systems.} LLM-MAS organize multiple
% role-specialized agents that communicate, plan, and act toward
% shared goals. Compared with a single agent, they can decompose
% long-horizon tasks, exploit diverse reasoning paths, cross-check
% intermediate results, and coordinate tool use~\cite{li2023camel,wu2023autogen,chen2023agentverse}.
% These properties have enabled role-playing collaboration, software
% engineering, open-ended problem solving, evaluation, and social
% simulation
% \cite{hong2024metagpt,qian2024chatdev,chan2023chateval}.
% Their modular structure also improves flexibility: agents can be
% added, replaced, or assigned different models without redesigning
% the whole workflow~\cite{du2024multiagentdebate,park2023generative}. Recent surveys therefore view LLM-MAS as a
% general architecture for scalable collective intelligence across
% heterogeneous tasks~\cite{guo2024survey}.

% \paragraph{LLM Multi-Agent Safeguarding.} Early safeguards often rely on prompt-based self-reflection, asking
% an agent to critique and revise suspicious outputs
% \cite{renze2024selfreflection}. This approach is lightweight, but it
% depends on the same model that may have produced the error and can
% become unreliable on long or tightly coupled trajectories.
% G-Safeguard instead applies a graph neural network to the agent
% utterance graph, detects anomalous agents, and prunes risky
% communication links \cite{wang2025gsafeguard}. Its supervised
% detector requires labeled malicious agents and may generalize poorly
% to unseen attacks \cite{miao2025blindguard}. Another common strategy
% introduces explicit reviewers. GuardAgent converts user-specified
% safety requirements into executable checks \cite{xiang2024guardagent},
% while TrinityGuard coordinates runtime monitor agents through a
% unified judge factory \cite{wang2026trinityguard}. PeerGuard uses
% mutual reasoning among agents to identify inconsistent peers
% \cite{fan2025peerguard}, and multi-agent debate relies on iterative
% cross-examination \cite{du2024multiagentdebate}. These methods improve
% oversight, but their effectiveness and cost vary substantially across
% complex tasks.

% However, existing safeguards either inspect interaction trajectories too coarsely, which delays detection, enlarges the rollback scope, and limits their effectiveness on long-horizon and tightly coupled tasks, or audit every interaction, which introduces substantial token and latency overhead. We therefore propose BRA-Audit, a budgeted audit point-placement framework that maintains effective trajectory auditing while reducing redundant checks. 

\paragraph{Attacks and Failures in LLM-MAS.}LLM-based multi-agent systems are vulnerable because agents exchange and reuse intermediate outputs, allowing local errors to escalate into system-level failures. Systematic analysis identifies failures in system design, inter-agent alignment, and task verification, revealing collaboration-specific risks beyond isolated agents~\cite{cemri2025why}. Adversaries can exploit the same dependencies: malicious instructions can self-propagate across agents~\cite{lee2024promptinfection}, manipulated knowledge can spread through conversations and persist in shared retrieval memories~\cite{ju2024flooding}, and communication attacks can intercept or modify inter-agent messages without compromising individual agents~\cite{he2025redteaming}. More severe control-flow hijacking can invoke unsafe tools, exfiltrate sensitive data, or execute malicious code~\cite{triedman2025maliciouscode}. Overall, communication, shared memory, and orchestration amplify benign errors and adversarial inputs, making early detection and containment essential for reliable LLM-MAS.


\paragraph{LLM Multi-Agent Safeguarding.}
Early guards use prompt-based self-reflection, but the reviewer may
share the generator's failure modes
~\cite{shinn2023reflexion,renze2024selfreflection}. G-Safeguard applies a graph neural
network to communication graphs to detect anomalous agents and prune
risky links, yet relies on labeled attacks and may generalize poorly to unseen threats
~\cite{wang2025gsafeguard,miao2025blindguard}. More reliable approaches
introduce LLM auditors. GuardAgent converts safety requirements
into executable checks, TrinityGuard coordinates runtime monitors,
PeerGuard uses mutual reasoning, and multi-agent debate performs
iterative cross-examination
~\cite{xiang2024guardagent,wang2026trinityguard,
fan2025peerguard,du2024multiagentdebate}. These methods improve
oversight, but their effectiveness and cost vary on complex tasks.

These safeguards methods either inspect trajectories too coarsely, which
weakens detection and enlarges the rollback scope, or audit every
interaction, which incurs substantial token and latency costs. We
therefore propose BRA-Audit, a budgeted audit-point placement framework
that preserves effective trajectory auditing while reducing redundant
checks.

\begin{figure*}[t]
    \centering
    \includegraphics[width=0.9\linewidth]{Figures/method.pdf}
    \caption{Overview of BRA-audit. }
    \label{fig:method}
\end{figure*}




% \section{System Model}
% \label{system model}
% This section formalizes the LLM-MAS execution model and the security threat model considered in this work.
% \subsection{MAS System}

% We consider a graph-structured LLM multi-agent system (LLM-MAS) with a set of agents
% \begin{equation}
% \mathcal{A}=\{a_1,\ldots,a_n\}.
% \end{equation}
% Each agent has a predefined role and performs a specific function in the task. Different roles may follow different interaction patterns, access different information sources or tools, and contribute to different stages of the overall execution.

% We model the runtime execution of the system as a dynamic dependency graph:
% \begin{equation}
% G_t=(V_t,E_t).
% \end{equation}
% Each node \(v\in V_t\) denotes a runtime object generated during execution, such as an agent message, an intermediate reasoning result, a tool request, a tool response, a retrieved piece of evidence, or an aggregation result. Each node is associated with its producing agent \(a(v)\), timestamp \(t(v)\), and content \(m(v)\). A directed edge \((u,v)\in E_t\) indicates that \(v\) depends on \(u\). Such a dependency can be induced by communication, evidence usage, tool-result usage, or aggregation. Thus, a path in \(G_t\) describes how information and potential errors propagate through the multi-agent execution.

% We model auditing as placing audit points on the execution graph. A audit point is a selected node \(s\in V_t\) inspected by an auditor. When \(s\) is selected, the auditor examines the audit context from the previous audit point to \(s\). This context contains the relevant messages, dependencies, and reasoning steps needed to judge whether the trajectory leading to \(s\) is safe. If the trajectory is judged safe, \(s\) becomes a trusted audit point. If it is judged unsafe, the system may block, sanitize, isolate, or roll back the affected trajectory.

% \subsection{Threat Model}

% We consider unsafe behavior that may arise during LLM-MAS execution. An agent may be directly compromised, or it may produce incorrect content due to hallucination, tool errors, poisoned knowledge bases, corrupted websites, or misleading retrieved documents. Although these causes are different, they have the same effect at the system level: an agent becomes a source of false, unsafe, or misleading information in the execution graph.

% We use the term \emph{malicious agent} as a unified abstraction for such cases. Let
% \begin{equation}
% \mathcal{M}\subseteq \mathcal{A},
% \end{equation}
% denote the set of malicious agents. For a node \(v\in V_t\), its content may be unsafe if it is produced by a malicious agent:
% \begin{equation}
% a(v)\in \mathcal{M}
% \Longrightarrow
% \Pr(y(v)=1)>0.
% \end{equation}
% Here, \emph{malicious} does not necessarily imply that the underlying LLM intentionally attacks the system. Instead, it means that the agent outputs corrupted information that can mislead downstream agents.

% A malicious agent may fabricate evidence, omit key constraints, provide an incorrect intermediate result, mislead an aggregation step, or trigger unsafe tool usage. Its effect can propagate along dependency edges. A corrupted message or intermediate result may be consumed by later agents, incorporated into subsequent reasoning, or used by an aggregator to form the final output. As a result, a local error can become a trajectory-level failure. In this work, we abstract these heterogeneous failure sources as malicious agents and study how runtime auditing can detect and contain their propagated effects over the execution graph.

% \begin{algorithm}[t]
% \caption{BRA: Exposure audit point Scheduling}
% \label{alg:bra_audit}
% \begin{algorithmic}[1]
% \REQUIRE Execution graph \(G_t=(V_t,E_t)\), candidate set \(Q_{t+1}\subseteq V_t\), audit ratio \(\rho\), trusted audit points \(\mathcal{H}_t\), last audit records \(\mathrm{LastAudit}_t(\cdot)\)
% \ENSURE Selected audit points \(S_{t+1}\), updated trusted audit points \(\mathcal{H}_{t+1}\)

% \STATE \(B_{t+1}\leftarrow \left\lfloor \rho |Q_{t+1}| \right\rfloor\)
% \STATE \(S_{t+1}\leftarrow \emptyset\)

% \FOR{each node \(u\in V_t\)}
%     \STATE \(\mathrm{AuditGap}_t(u)\leftarrow t-\mathrm{LastAudit}_t(a(u))+1\)
%     \STATE \(\mathrm{TopoImpact}_t(u)\leftarrow 1+|\mathrm{Desc}_t(u)|\)
%     \STATE \(w_t(u)\leftarrow \mathrm{AuditGap}_t(u)\cdot \mathrm{TopoImpact}_t(u)\)
% \ENDFOR

% \WHILE{\(|S_{t+1}|<B_{t+1}\) and \(Q_{t+1}\setminus S_{t+1}\neq \emptyset\)}
%     \FOR{each candidate \(v\in Q_{t+1}\setminus S_{t+1}\)}
%         \STATE \(\Delta(v\mid S_{t+1})\leftarrow J(\mathcal{H}_t\cup S_{t+1})-J(\mathcal{H}_t\cup S_{t+1}\cup\{v\})\)
%     \ENDFOR
%     \STATE \(v^\star\leftarrow \arg\max_{v\in Q_{t+1}\setminus S_{t+1}}\Delta(v\mid S_{t+1})\)
%     \STATE \(S_{t+1}\leftarrow S_{t+1}\cup\{v^\star\}\)
% \ENDWHILE

% \STATE \(\mathcal{H}_{t+1}\leftarrow \mathcal{H}_t\)

% \FOR{each audit point \(s\in S_{t+1}\)}
%     \STATE \(\mathcal{T}(s)\leftarrow \operatorname{Path}(\operatorname{Prev}(s\mid \mathcal{H}_{t+1}\cup S_{t+1}),s)\)
%     \STATE \(r\leftarrow \operatorname{Auditor}(\mathcal{T}(s))\)
%     \IF{\(r=\mathrm{safe}\)}
%         \STATE \(\mathcal{H}_{t+1}\leftarrow \mathcal{H}_{t+1}\cup\{s\}\)
%         \STATE update \(\mathrm{LastAudit}_{t+1}(a(s))\)
%     \ELSE
%         \STATE rollback the work products affected by \(\mathcal{T}(s)\)
%         \STATE mark nodes in \(\mathcal{T}(s)\) as unsafe or suspicious
%     \ENDIF
% \ENDFOR

% \RETURN \(S_{t+1}, \mathcal{H}_{t+1}\)
% \end{algorithmic}
% \end{algorithm}

\begin{algorithm}[t]
\caption{BRA: Exposure audit-point Scheduling}
\label{alg:bra_audit}
\begin{algorithmic}[1]
\REQUIRE Execution graph \(G_t=(V_t,E_t)\), candidate set
\(Q_{t+1}\subseteq V_t\), audit ratio \(\rho\), and last audit
records \(\mathrm{LastAudit}_t(\cdot)\)
\ENSURE Selected audit points \(S_{t+1}\) and updated audit
records \(\mathrm{LastAudit}_{t+1}(\cdot)\)

\STATE \(B_{t+1}\leftarrow
\left\lfloor \rho |Q_{t+1}| \right\rfloor\)
\STATE \(S_{t+1}\leftarrow \emptyset\)

\FOR{each node \(u\in V_t\)}
    \STATE \(\mathrm{AuditGap}_t(u)
    \leftarrow t-\mathrm{LastAudit}_t(a(u))+1\)
    \STATE \(\mathrm{TopoImpact}_t(u)
    \leftarrow 1+|\mathrm{Desc}_t(u)|\)
    \STATE \(w_t(u)\leftarrow
    \mathrm{AuditGap}_t(u)\cdot
    \mathrm{TopoImpact}_t(u)\)
\ENDFOR

\WHILE{\(|S_{t+1}|<B_{t+1}\) and
\(Q_{t+1}\setminus S_{t+1}\neq\emptyset\)}
    \FOR{each \(v\in Q_{t+1}\setminus S_{t+1}\)}
        \STATE \(\Delta(v\mid S_{t+1})
        \leftarrow J(S_{t+1})
        -J(S_{t+1}\cup\{v\})\)
    \ENDFOR
    \STATE \(v^\star\leftarrow
    \arg\max_{v\in Q_{t+1}\setminus S_{t+1}}
    \Delta(v\mid S_{t+1})\)
    \STATE \(S_{t+1}\leftarrow
    S_{t+1}\cup\{v^\star\}\)
\ENDWHILE

\STATE deploy \(S_{t+1}\) before round-\((t+1)\) execution

\FOR{each audit point \(s\in S_{t+1}\) in causal order}
    \STATE \(\mathcal{T}(s)\leftarrow
    \operatorname{Region}(s\mid S_{t+1})\)
    \STATE \(r\leftarrow
    \operatorname{Auditor}(\mathcal{T}(s))\)

    \IF{\(r=\mathrm{safe}\)}
        \STATE mark \(s\) as a trusted audit point for
        downstream execution in this round
    \ELSE
        \STATE block the affected downstream information flow
        \STATE rollback and recover the work products
        affected by \(\mathcal{T}(s)\)
        \STATE re-execute and re-audit \(\mathcal{T}(s)\)
        before continuing
    \ENDIF

    \STATE update
    \(\mathrm{LastAudit}_{t+1}(a(s))\)
\ENDFOR

\RETURN \(S_{t+1},
\mathrm{LastAudit}_{t+1}(\cdot)\)
\end{algorithmic}
\end{algorithm}

\section{System Model}
\label{system model}

This section defines the LLM-MAS execution model and the threat model considered in this work.

\subsection{MAS System}

We consider a graph-structured LLM multi-agent system (LLM-MAS) with a set of agents
\begin{equation}
\mathcal{A}=\{a_1,\ldots,a_n\}.
\end{equation}
Each agent has a predefined role and contributes to the task through its own interaction pattern, information sources, and tools. We represent runtime execution as a dynamic dependency graph:
\begin{equation}
G_t=(V_t,E_t).
\end{equation}
Each node \(v\in V_t\) represents a runtime object, such as an agent message, intermediate result, tool interaction, retrieved evidence, or aggregation result. It is associated with its producing agent \(a(v)\), timestamp \(t(v)\), and content \(m(v)\). A directed edge \((u,v)\in E_t\) indicates that \(v\) depends on \(u\) through communication, evidence usage, tool execution, or aggregation. A path in \(G_t\) therefore captures the propagation of information and potential errors.

Auditing is modeled as audit-point placement on the execution graph. An audit point is a selected node \(s\in V_t\) inspected by an auditor using the relevant context between the previous audit point and \(s\). If the trajectory is safe, \(s\) becomes a trusted audit point in this round. Otherwise, the system may block, sanitize, isolate, or roll back the affected trajectory.

\subsection{Threat Model}

Unsafe behavior may arise from compromised agents, hallucinations, tool errors, poisoned knowledge bases, corrupted websites, or misleading retrieved documents. Despite their different causes, these failures introduce false, unsafe, or misleading information into the execution graph.

We use \emph{malicious agents} as a unified abstraction for these cases. Let
\begin{equation}
\mathcal{M}\subseteq \mathcal{A},
\end{equation}
denote the set of malicious agents. For a node \(v\in V_t\), its content may be unsafe if it is produced by a malicious agent:
\begin{equation}
a(v)\in \mathcal{M}
\Longrightarrow
\Pr(y(v)=1)>0.
\end{equation}
Here, \(y(v)=1\) denotes unsafe content. The term \emph{malicious} does not require intentional behavior; it only indicates that an agent produces corrupted information that may mislead downstream agents.

A malicious agent may fabricate evidence, omit constraints, produce incorrect intermediate results, distort aggregation, or trigger unsafe tool usage. These outputs can propagate through dependency edges and affect later reasoning or the final result. We therefore study how runtime auditing detects and contains such propagated failures over the execution graph.


\section{Method}
This section formulates budgeted runtime auditing as an optimization problem and develops the mathematical model and scheduling algorithm used by BRA-Audit (Fig. ~\ref{fig:method} shows overview of BRA-audit).


\subsection{Problem Formulation and Optimization Objective}

% LLM-based auditors can inspect intermediate MAS outputs and identify hallucinated, inconsistent, or unsafe reasoning
% steps~\citep{zheng2023judgingllmasajudgemtbenchchatbot,
% gu2025surveyllmasajudge}. As discussed in the introduction, final-only auditing requires processing a long execution trajectory and may cause large rollback, whereas auditing every interaction requires frequent auditor calls. We therefore study how to allocate a limited audit budget during MAS execution.

% At the beginning of each audit round \(t\), BRA-Audit selects a new set of audit points before the corresponding information flow is executed. The selected nodes act as audit triggers. When execution reaches a audit point, the auditor inspects the dependency region accumulated since the preceding audit points on its incoming dependency paths. A successful audit establishes the inspected node as a trusted audit point for subsequent downstream execution in the same round. The resulting audit point-bounded regions determine both the context that the auditor must inspect and the execution scope that may need to be rolled back.

% The locations of audit points are also important. If a node influences many downstream computations, leaving it unchecked can introduce persistent exposure into later audit contexts. Moreover, audit point decisions are coupled because placing one audit point changes the unchecked dependency regions and the value of subsequent audit points.

% These observations motivate us to formulate budgeted auditing as a audit point-placement problem. Given a audit point set
% \begin{equation}
% S_t \subseteq V_t,
% \qquad
% |S_t| \leq B_t,
% \end{equation}
% we select the audit points before execution so that the resulting audit point-bounded audit regions do not contain overly long reasoning trajectories. At the same time, the scheduler considers the topological importance of nodes and prioritizes agents or regions that have not been audited for many rounds. In a dependency graph, a corrupted node with broad downstream reachability can impose higher auditing and rollback costs on subsequent computation, while a long audit gap indicates stale trust evidence for the corresponding agent or region and makes an earlier audit point more valuable.

% We next formalize this intuition as a cumulative-exposure audit point-placement objective:
% \begin{equation}
% \min_{S_t\subseteq V_t}
% J(S_t)
% \quad
% \mathrm{s.t.}
% \quad
% |S_t|\leq B_t,
% \end{equation}
% where \(J(S_t)\) measures the total unchecked exposure induced by the audit point placement \(S_t\). We define this objective and its components in the following subsection.
LLM-based auditors can inspect intermediate MAS outputs and identify hallucinated, inconsistent, or unsafe reasoning
steps~\citep{zheng2023judgingllmasajudgemtbenchchatbot,
gu2025surveyllmasajudge}. As discussed in the introduction, final-only auditing requires processing a long execution trajectory and may cause large rollback, whereas auditing every interaction requires frequent auditor calls. We therefore study how to allocate a limited audit budget during MAS execution.

A successful audit verifies the audit point-bounded dependency
region reaching the inspected node and establishes the node as
a trusted audit point in this round for subsequent execution. The size of this
region determines both the context inspected by the auditor and
the execution scope that may need to be rolled back. Audit-point
placement is therefore important: leaving a node with broad
downstream reachability unchecked can repeatedly carry its
exposure into later audit contexts, while a long audit gap
indicates stale trust evidence for the corresponding agent or
region. Moreover, audit point decisions are coupled because each
new audit point reshapes downstream audit regions and changes the
value of subsequent audit points.

Given a candidate set \(Q_t\subseteq V_t\) and an audit-call
budget \(B_t\), BRA-Audit selects before execution an audit point
set
\begin{equation}
S_t\subseteq Q_t,
\qquad
|S_t|\leq B_t.
\end{equation}
Each selected audit point triggers one auditor call over its
audit-point bounded dependency region. The scheduler therefore
seeks placements that keep audit regions compact while
prioritizing influential nodes and agents that have remained
unaudited for many rounds.

We formulate audit point placement as minimizing cumulative
unchecked exposure:
\begin{equation}
\min_{S_t\subseteq Q_t} J(S_t)
\quad
\mathrm{s.t.}
\quad
|S_t|\leq B_t,
\end{equation}
where \(J(S_t)\) measures the total cumulative exposure
remaining under audit-point placement \(S_t\). By preventing
unchecked information from being repeatedly inherited by
downstream contexts, this objective jointly reduces audit
context size and delayed-detection exposure. Here, \(B_t\)
controls the number of auditor calls, while total token
cost is evaluated end to end. We define \(J(S_t)\) and
its components in the next subsection.

% A successful audit verifies the trajectory leading to the inspected node and establishes a trusted audit point in the execution graph. The distance between two such boundaries determines both the context that the auditor must inspect and the execution scope that may need to be rolled back. Their locations are also important. If a node influences many downstream computations, leaving it unchecked can introduce persistent exposure into later audit contexts. Moreover, audit decisions are coupled because establishing one trusted audit point changes the unchecked contexts and the value of subsequent audits.

% These observations motivate us to formulate budgeted auditing as a audit point-placement problem. Given a audit point set\begin{equation} 
% S_t \subseteq V_t, \qquad |S_t| \leq B_t ,
% \end{equation}
% we need to insert them into the interaction graph so that neighboring audit points are not separated by overly long reasoning trajectories. At the same time, the scheduler should consider the topological importance of nodes and should prioritize agents or regions that have not been audited for many rounds. This is because, in a dependency graph, a corrupted node with broad downstream reachability can impose higher auditing and rollback costs on subsequent computation, while a long audit gap indicates stale trust evidence for the corresponding agent or region and therefore makes earlier audit points more desirable. We next formalize this intuition as a cumulative-exposure audit point placement and minimizing a cumulative-exposure objective: 
% \begin{equation} \min_{S_t\subseteq V_t} J(S_t) \quad \mathrm{s.t.} \quad |S_t|\leq B_t , 
% \end{equation} 
% where \(J(S_t)\) measures the total unaudited exposure induced by the audit point set \(S_t\). We define this objective and its components in the following subsection.


% \subsection{Analysis and Solution}

% We now analyze the cumulative-exposure objective and derive the audit point scheduling rule used by BRA-Audit. Algorithm~\ref{alg:bra_audit} summarizes the our main idea of BRA-Audit. The key observation is that an unchecked node does not only affect its own audit context. If it is not separated by a audit point, its exposure continues to appear in the audit contexts of downstream nodes. Thus, audit scheduling should not be treated as an independent top-\(k\) node selection problem. Instead, it is a audit point placement problem that aims to cut off the cumulative propagation of unchecked exposure.

% We first define the exposure weight of each node. For a node \(u\in V_t\), let \(a(u)\) denote its producing agent. We define the audit gap of \(u\) as
% \begin{equation}
% \mathrm{AuditGap}_t(u)
% =
% t-\mathrm{LastAudit}_t(a(u))+1,
% \end{equation}
% where \(\mathrm{LastAudit}_t(a(u))\) is the most recent round before \(t\) in which agent \(a(u)\) was audited. A larger audit gap means that the producing agent has remained unaudited for more rounds.

% We also define the topological impact of \(u\) by its downstream reachability in the dependency graph:
% \begin{equation}
% \mathrm{TopoImpact}_t(u)
% =
% 1+\left|\mathrm{Desc}_t(u)\right|,
% \end{equation}
% where \(\mathrm{Desc}_t(u)\) is the set of downstream nodes that depend on \(u\) in \(G_t\). This term captures the potential propagation scope of an unsafe output. Moreover, topologically important nodes tend to play more critical roles in task execution and should therefore receive higher audit priority. The exposure weight of \(u\) is then defined as
% \begin{equation}
% w_t(u)
% =
% \mathrm{AuditGap}_t(u)\cdot \mathrm{TopoImpact}_t(u).
% \end{equation}
% This weight uses only runtime logs and graph structure. It reflects two observable factors: how long the producing agent has not been audited, and how broadly the node can affect downstream computation.

% Given a audit point set \(S_t\), let \(\mathrm{Ctx}(v\mid S_t)\) denote the unchecked audit context of node \(v\). This context contains the nodes from the nearest preceding audit point to \(v\) that have not been cut off by another audit point. The cumulative exposure of \(v\) is
% \begin{equation}
% C(v\mid S_t)
% =
% \sum_{u\in \mathrm{Ctx}(v\mid S_t)} w_t(u).
% \end{equation}
% The total cumulative exposure is
% \begin{equation}
% J(S_t)
% =
% \sum_{v\in V_t} C(v\mid S_t)
% =
% \sum_{v\in V_t}
% \sum_{u\in \mathrm{Ctx}(v\mid S_t)}
% w_t(u).
% \end{equation}
% This objective measures how much unchecked exposure remains in all audit contexts. A high-exposure node that is not audit pointed can repeatedly appear in later contexts. This creates a persistent audit burden and increases the cost of delayed detection.

% Directly minimizing \(J(S_t)\) is difficult in a dynamic execution graph. The number of possible audit point sets grows quickly with \(|V_t|\). In addition, audit point decisions are coupled. Adding a audit point changes the audit contexts of downstream nodes, so the value of later audit points depends on earlier choices. BRA-Audit therefore uses an incremental marginal-gain strategy.

% Given the current audit point set \(S_t\), the marginal gain of adding a candidate audit point \(v\in V_t\setminus S_t\) is defined as
% \begin{equation}
% \Delta(v\mid S_t)
% =
% J(S_t)-J(S_t\cup\{v\}).
% \end{equation}
% This quantity measures how much total cumulative exposure is reduced if \(v\) is selected as a audit point. A useful audit point should cut off a large amount of accumulated exposure and prevent it from entering many downstream audit contexts.

% To make this effect explicit, consider a candidate audit point \(v\). Let \(\mathrm{Pre}(v\mid S_t)\) denote the unchecked prefix from the latest preceding audit point to \(v\), including \(v\). Let \(\mathrm{Down}(v\mid S_t)\) denote the downstream audit contexts that would otherwise continue to contain this prefix. Then the reduction caused by selecting \(v\) can be written as
% \begin{equation}
% \Delta(v\mid S_t)
% =
% \sum_{z\in \mathrm{Down}(v\mid S_t)}
% \sum_{u\in \mathrm{Pre}(v\mid S_t)}
% w_t(u).
% \end{equation}
% This expression states that the value of \(v\) depends on two factors: the accumulated exposure before \(v\), and the number of downstream contexts that would continue to carry this exposure. Therefore, a audit point is valuable when it cuts off a high-exposure prefix that would otherwise affect many later nodes. A concrete calculation for a chain-structured trajectory is provided in Technical Appendix section example.

% BRA-Audit selects audit points according to this marginal gain. At each step, it selects the candidate that gives the largest reduction in cumulative exposure:
% \begin{equation}
% v^\star
% =
% \arg\max_{v\in V_t\setminus S_t}
% \Delta(v\mid S_t).
% \end{equation}
% The selected node is added to the audit point set, and the affected audit contexts are updated. This process is repeated until the audit budget is used. The resulting audit point set provides a practical approximation to the cumulative-exposure minimization problem (We further justify the design and provide a concrete example in the technical supplement).


\subsection{Analysis and Solution}

We now analyze the cumulative-exposure objective and derive the audit-point scheduling rule used by BRA-Audit. Algorithm~\ref{alg:bra_audit} summarizes the main procedure of BRA-Audit. The key observation is that an unchecked node does not only affect its own audit context. Its exposure can continue to appear in downstream dependency contexts until the corresponding information flow reaches an audit point. Audit scheduling therefore considers how an audit point jointly reduces exposure in multiple downstream contexts.

Since runtime objects are ordered by their generation times and dependency edges follow causal information flow, repeated agent interactions are unrolled across rounds, and the runtime dependency graph \(G_t\) is represented as a directed acyclic graph.

% We first define the exposure weight of each node. For a node \(u\in V_t\), let \(a(u)\) denote its producing agent. We define the audit gap of \(u\) as
% \begin{equation}
% \mathrm{AuditGap}_t(u)
% =
% t-\mathrm{LastAudit}_t(a(u))+1,
% \end{equation}
% where \(\mathrm{LastAudit}_t(a(u))\) is the most recent round before \(t\) in which a audit point produced by agent \(a(u)\) was audited. A larger audit gap means that the producing agent has remained unaudited for more rounds.

% We also define the topological impact of \(u\) by its downstream reachability in the dependency graph:
% \begin{equation}
% \mathrm{TopoImpact}_t(u)
% =
% 1+\left|\mathrm{Desc}_t(u)\right|,
% \end{equation}
% where \(\mathrm{Desc}_t(u)\) is the set of downstream nodes that depend on \(u\) in \(G_t\). This term captures the potential propagation scope of an unsafe output. Topologically important nodes tend to play more critical roles in task execution and therefore receive higher audit priority. The exposure weight of \(u\) is defined as
% \begin{equation}
% w_t(u)
% =
% \mathrm{AuditGap}_t(u)\cdot
% \mathrm{TopoImpact}_t(u).
% \end{equation}
% This weight uses only runtime logs and graph structure. It reflects two observable factors: how long the producing agent has not been audited and how broadly the node can affect downstream computation.

% For each node \(u\in V_t\), let \(a(u)\) denote its producing
% agent. We define its audit gap, topological impact, and
% exposure weight as
% \begin{equation}
% \begin{aligned}
% \mathrm{AuditGap}_t(u)&=t-\mathrm{LastAudit}_t(a(u))+1,\\
% \mathrm{TopoImpact}_t(u)&=1+\left|\mathrm{Desc}_t(u)\right|,\\
% w_t(u)&=\mathrm{AuditGap}_t(u)\,
%   \mathrm{TopoImpact}_t(u).
% \end{aligned}
% \end{equation}
% Here, \(\mathrm{LastAudit}_t(a(u))\) is the latest round before
% \(t\) in which an output of agent \(a(u)\) was audited, and
% \(\mathrm{Desc}_t(u)\) is the set of nodes downstream of \(u\).
% Thus, \(\mathrm{AuditGap}_t(u)\) captures stale audit coverage,
% while \(\mathrm{TopoImpact}_t(u)\) estimates the potential
% propagation scope of an unsafe output. Nodes with broad
% downstream reachability also tend to occupy more important
% topological positions and play more critical roles in
% information aggregation, propagation, or decision making.
% Their product prioritizes nodes that are both long unaudited
% and influential, using only runtime logs and graph structure.

For each node \(u\in V_t\), let \(a(u)\) denote its producing agent, and let \(\tau_t(a)\) be the latest completed audit round of agent \(a\) available at scheduling time \(t\). We define 
\begin{equation} \begin{aligned} \mathrm{Audit Gap}(u) &= (t+1)-\tau_t(a(u)),\\ \mathrm{TopoImpact}_t(u) &= \left|\{u\}\cup\mathrm{Desc}_t(u)\right|,\\ w_{t+1\mid t}(u) &= \mathrm{AuditGap}(u)\, \mathrm{TopoImpact}_t(u), 
\end{aligned} 
\end{equation} 

where \(\mathrm{Desc}_t(u)\) denotes the strict descendants of \(u\).
We consider verification age because risk can accumulate while an
agent's outputs remain unchecked, and downstream reach because errors
at influential nodes can affect more subsequent computations and nodes with broad downstream reachability also tend to play more critical roles in the task. Including
\(u\) itself ensures that even sink nodes retain nonzero local impact.
Accordingly, the first term measures the verification age when the
scheduled audit is executed, while the second measures the closed
downstream region potentially affected by \(u\). Their product serves
as a temporal--topological exposure measure, prioritizing nodes with
stale verification and broad downstream reach using runtime audit
records and graph structure.



Let \(\mathrm{Anc}_t(v)\) denote the ancestors of \(v\).
We write \(u\leadsto_{S_t}v\) if there exists a directed path
\(p:u\leadsto v\) such that
\((V(p)\setminus\{v\})\cap S_t=\varnothing\).
Given an audit point set \(S_t\), the unchecked audit context
reaching \(v\) is
\begin{equation}
\mathrm{Ctx}(v\mid S_t)
=
\left\{
u\in\mathrm{Anc}_t(v)\cup\{v\}
:
u\leadsto_{S_t}v
\right\}.
\end{equation}
Equivalently, it is obtained by traversing backward from \(v\)
and stopping each branch at the first selected audit point.
This definition preserves unresolved branches in a DAG: if
any path from \(u\) to \(v\) bypasses the selected audit points,
\(u\) remains in the audit context. Thus, an audit point shortens
only the downstream contexts whose information flow passes
through it.

All audit points are selected before execution and triggered
in causal order. Once an audit point passes the audit, its
output becomes trusted input for downstream execution; an
unsafe region is blocked and recovered before execution
continues.

The cumulative exposure of \(v\) and the total exposure are
\begin{equation}
\begin{aligned}
C(v\mid S_t)
&=
\sum_{u\in\mathrm{Ctx}(v\mid S_t)}w_t(u),\\
J(S_t)
&=
\sum_{v\in V_t}C(v\mid S_t).
\end{aligned}
\end{equation}
A high-exposure node can repeatedly contribute to downstream
contexts until bounded by an audit point, increasing the audit
burden and the cost of delayed detection. By summing over all
nodes, \(J(S_t)\) captures both the size of local audit contexts
and the repeated inheritance of unchecked exposure by
downstream computation.

Directly minimizing \(J(S_t)\) is combinatorial, and
audit point decisions are coupled because each selected
audit point changes downstream contexts and the value of
later selections. BRA-Audit therefore greedily selects
audit points by marginal gain.

\begin{table*}[t]
\centering
\caption{Performance and token cost on AgentsNet.
Token consumption is reported in thousands (K).}
\label{tab:agentsnet_main}

\small
\setlength{\tabcolsep}{6pt}
\renewcommand{\arraystretch}{1.10}

\resizebox{0.9\linewidth}{!}{%
\begin{tabular}{@{}l*{12}{c}@{}}
\toprule
& \multicolumn{2}{c}{\textbf{Consensus}}
& \multicolumn{2}{c}{\textbf{Leader}}
& \multicolumn{2}{c}{\textbf{Matching}}
& \multicolumn{2}{c}{\textbf{Coloring}}
& \multicolumn{2}{c}{\textbf{VertexCover}}
& \multicolumn{2}{c}{\textbf{Overall}} \\
\cmidrule(lr){2-3}
\cmidrule(lr){4-5}
\cmidrule(lr){6-7}
\cmidrule(lr){8-9}
\cmidrule(lr){10-11}
\cmidrule(lr){12-13}

\textbf{Method}
& \textbf{Score} & \textbf{Tok.}
& \textbf{Score} & \textbf{Tok.}
& \textbf{Score} & \textbf{Tok.}
& \textbf{Score} & \textbf{Tok.}
& \textbf{Score} & \textbf{Tok.}
& \textbf{Score} & \textbf{Tok.} \\
\midrule

Clean
& 0.9630 & 28.1
& 0.5926 & 29.9
& 0.6944 & 25.7
& 0.9967 & 31.1
& 0.8342 & 18.6
& 0.8162 & 26.7 \\

Attack
& 0.0741 & 42.2
& 0.1481 & 43.1
& 0.4815 & 35.7
& 0.8324 & 47.5
& 0.6676 & 29.0
& 0.4407 & 39.5 \\

\midrule

PeerGuard
& 0.9630 & 69.3
& 0.5556 & 67.5
& 0.8056 & 68.1
& 0.9716 & 70.4
& 0.8852 & 53.6
& 0.8362 & 65.8 \\

AuditAgent
& 0.8148 & 48.6
& 0.4444 & 50.2
& 0.6250 & 49.1
& 0.9013 & 53.9
& 0.7854 & 40.3
& 0.7142 & 48.4 \\

GuardAgent
& 0.4074 & 44.7
& 0.3704 & 46.5
& 0.5370 & 39.3
& 0.8548 & 50.2
& 0.7014 & 33.4
& 0.5742 & 42.8 \\

TrinityGuard
& 0.8889 & 58.0
& 0.4815 & 55.6
& 0.7385 & 51.8
& 0.8892 & 64.8
& 0.8156 & 44.9
& 0.7627 & 55.0 \\

Multi-Agent Debate
& 1.0000 & 75.4
& 0.6296 & 74.5
& 0.8342 & 77.3
& 0.9967 & 81.6
& 0.8675 & 71.5
& 0.8656 & 76.1 \\

\midrule

\textbf{BRA-Audit}
& 0.9630 & 57.6
& 0.5185 & 57.0
& 0.7917 & 52.4
& 0.9967 & 61.7
& 0.8342 & 43.8
& 0.8208 & 54.5 \\

\bottomrule
\end{tabular}%
}
\end{table*}

For a candidate \(v\in V_t\setminus S_t\), let
\(\mathrm{Cut}_v(z\mid S_t)\) denote the exposure removed
from the context of \(z\):
\begin{equation}
\begin{aligned}
\mathrm{Cut}_v(z\mid S_t)
&=
\mathrm{Ctx}(z\mid S_t)\\
&\quad\setminus
\mathrm{Ctx}(z\mid S_t\cup\{v\}).
\end{aligned}
\end{equation}
The marginal gain of \(v\) is
\begin{equation}
\begin{aligned}
\Delta(v\mid S_t)
&=
J(S_t)-J(S_t\cup\{v\})\\
&=
\sum_{z\in V_t}
\sum_{u\in\mathrm{Cut}_v(z\mid S_t)}
w_t(u).
\end{aligned}
\end{equation}
This quantity measures the exposure removed from all
dependency contexts by adding \(v\). It is large when \(v\)
separates a high-exposure region from many downstream
contexts. Since marginal gains are recomputed after each
selection, the scheduler accounts for overlapping effects
among audit points rather than scoring them independently.
For a chain, the expression reduces to the accumulated prefix
exposure multiplied by the number of affected downstream
contexts, as illustrated in the technical appendix.





At each step, BRA-Audit selects
\begin{equation}
v^\star
=
\arg\max_{v\in V_t\setminus S_t}
\Delta(v\mid S_t),
\end{equation}
adds \(v^\star\) to \(S_t\), and updates the affected contexts
until the audit budget is exhausted.


\begin{table*}[t]
\centering
\caption{Performance and token cost on BBH and Multi-agent-bench-rasearch.}
\label{tab:BBH}
\resizebox{0.9\linewidth}{!}{%
\begin{tabular}{cccccccccc}
\hline
                            & \multicolumn{3}{c}{\textbf{BBH}}                  & \multicolumn{6}{c}{\textbf{MultiAgentBench-Research}}                                                                 \\ \hline
\textbf{Method}             & \textbf{Acc} & \textbf{recovery} & \textbf{token} & \textbf{innovation} & \textbf{safety} & \textbf{feasibility} & \textbf{Score} & \textbf{recovery} & \textbf{token(K)} \\ \hline
\textbf{clean}              & 0.9200       & --                & 5355           & 5.00                & 4.62            & 3.83                 & 0.897          & --                & 23.0              \\
\textbf{attack}             & 0.3933       & --                & 6778           & 2.28                & 1.88            & 1.50                 & 0.377          & --                & 29.1              \\
\textbf{PeerGuard}          & 0.8867       & 0.4933            & 10136          & 4.96                & 4.80            & 3.86                 & 0.908          & 0.531             & 47.3              \\
\textbf{Auditagent}         & 0.8067       & 0.4133            & 7459           & 3.88                & 3.92            & 3.14                 & 0.730          & 0.352             & 32.4              \\
\textbf{Guardagent}         & 0.5200       & 0.1267            & 7318           & 3.20                & 2.74            & 2.88                 & 0.588          & 0.211             & 30.8              \\
\textbf{Trinityguard}       & 0.8600       & 0.4667            & 8324           & 4.32                & 4.16            & 3.64                 & 0.808          & 0.431             & 37.9              \\
\textbf{Multi-Agent Debate} & 0.9333       & 0.54              & 13763          & 4.98                & 4.86            & 4.02                 & 0.925          & 0.548             & 60.5              \\ \hline
\textbf{BRA}                & 0.8600       & 0.4667            & 8175           & 5.00                & 4.74            & 3.78                 & 0.902          & 0.525             & 38.5              \\ \hline
\end{tabular}
}
\end{table*}



\subsection{Online audit-point Placement and Audit Execution}

BRA-Audit operates in an online manner. At the end of each MAS interaction round \(t\), the system updates the dependency graph \(G_t=(V_t,E_t)\). Before the next round starts, the scheduler uses this graph to place a limited number of audit points. Let \(Q_{t+1}\subseteq V_t\) denote the candidate nodes for the next audit round. Given an audit ratio \(\rho\in(0,1]\), which corresponds to auditing \(k\%\) of the candidates, the audit budget and audit point set are defined as
\begin{equation}
B_{t+1}=\left\lfloor \rho |Q_{t+1}| \right\rfloor,
\qquad
S_{t+1}\approx
\arg\min_{S\subseteq Q_{t+1},\, |S|\leq B_{t+1}} J(S).
\end{equation}
Here, \(J(S)\) is the cumulative unaudited exposure objective defined above. In practice, BRA-Audit selects audit points that provide the largest reduction of \(J(S)\) until the budget is used. In the final round, the agent producing the final result is mandatorily selected and counted toward the budget.

% After the audit point set \(S_{t+1}\) is selected, each audit point is assigned to an auditor. For a audit point \(s\in S_{t+1}\), the auditor inspects the local audit trajectory from the nearest preceding audit point to \(s\):
% \begin{equation}
% \mathcal{T}(s\mid S_{t+1})
% =
% \operatorname{Path}\big(\operatorname{Prev}(s\mid S_{t+1}),\,s\big).
% \end{equation}

After the audit-point set \(S_{t+1}\) is selected, each
audit point is assigned to an auditor. For an audit point
\(s\in S_{t+1}\), the auditor inspects the audit-point bounded
dependency region accumulated along all incoming branches:
\begin{equation}
\mathcal{T}(s\mid S_{t+1})
=
\operatorname{Region}(s\mid S_{t+1}).
\end{equation}

Here, \(\operatorname{Region}(s\mid S_{t+1})\) is obtained by
traversing backward from \(s\) and stopping each branch at
the first upstream audit point in \(S_{t+1}\). The audit point
outputs are treated as trusted inputs, while the subsequent
inputs, outputs, messages, and dependency steps form the local
audit region.

If the auditor judges \(\mathcal{T}(s\mid S_{t+1})\) as safe,
\(s\) becomes a trusted audit point; otherwise, the system
rolls back the work products affected by this region. By
partitioning execution into audit-point bounded regions,
BRA-Audit shortens audit contexts and localizes rollback.
Because the auditor examines the entire region rather than
only its endpoint, an unsafe agent can still be detected even
when it is not selected as an audit point, provided that its
inputs or outputs appear in the region of an downstream
audit point.
% Here, \(\operatorname{Region}(s\mid S_{t+1})\) is obtained by
% traversing the dependency graph backward from \(s\) and
% stopping each branch at the first upstream audit point in
% \(S_{t+1}\). The audit point-audit point outputs are included as
% trusted inputs, while the nodes accumulated after these
% boundaries form the local audit region.

% This trajectory contains the relevant inputs, outputs, messages, and dependency steps generated by the agents along the path. If the auditor judges \(\mathcal{T}(s\mid S_{t+1})\) as safe, \(s\) is recorded as a trusted audit point. While unsafe, the system rolls back the work products affected by this audit trajectory.

% This procedure makes auditing local rather than global. By splitting execution into audit point-bounded audit segments, BRA-Audit shortens audit contexts and localizes rollback relative to long-trace auditing. Moreover, the auditor checks the whole trajectory between audit points rather than only the endpoint. Thus, even if an unsafe agent is not itself selected as a audit point in one round, its inputs and outputs may still be exposed when they lie on the audit trajectory of a downstream audit point.

% \begin{table*}[t]
% \centering
% \caption{Backbone Model Ablation.}
% \label{tab:model}
% \begin{tabular}{lcccccccccccc}
% \hline
%                           & \multicolumn{2}{c}{\textbf{Consensus}}                                     & \multicolumn{2}{c}{\textbf{Leader}}                                        & \multicolumn{2}{c}{\textbf{Matching}}                                      & \multicolumn{2}{c}{\textbf{Coloring}}                                      & \multicolumn{2}{c}{\textbf{VertexCover}}                                   & \multicolumn{2}{c}{\textbf{Overall}}                                       \\ \hline
% \textbf{Model}            & \multicolumn{1}{l}{\textbf{Score}} & \multicolumn{1}{l}{\textbf{Token(K)}} & \multicolumn{1}{l}{\textbf{Score}} & \multicolumn{1}{l}{\textbf{Token(K)}} & \multicolumn{1}{l}{\textbf{Score}} & \multicolumn{1}{l}{\textbf{Token(K)}} & \multicolumn{1}{l}{\textbf{Score}} & \multicolumn{1}{l}{\textbf{Token(K)}} & \multicolumn{1}{l}{\textbf{Score}} & \multicolumn{1}{l}{\textbf{Token(K)}} & \multicolumn{1}{l}{\textbf{Score}} & \multicolumn{1}{l}{\textbf{Token(K)}} \\ \hline
% \textbf{Qwen}             & 0.9630                             & 57.6                                  & 0.5185                             & 57.0                                  & 0.7917                             & 52.4                                  & 0.9967                             & 61.7                                  & 0.8342                             & 43.8                                  & 0.8208                             & 54.5                                  \\
% \textbf{DeepseekV4-flash} & 0.9630                             & 51.0                                  & 0.5556                             & 72.9                                  & 0.6435                             & 64.2                                  & 0.9463                             & 51.3                                  & 0.8250                             & 44.6                                  & 0.7867                             & 56.8                                  \\
% \textbf{Minmax2.7}        & 0.9259                             & 55.8                                  & 0.5185                             & 63.4                                  & 0.7592                             & 51.3                                  & 0.9177                             & 55.4                                  & 0.7935                             & 35.6                                  & 0.783                              & 52.3                                  \\
% \textbf{Glm5.1}           & 0.9259                             & 62.9                                  & 0.5926                             & 63.2                                  & 0.7314                             & 56.8                                  & 0.9446                             & 61.0                                  & 0.8250                             & 53.1                                  & 0.8039                             & 59.4                                  \\
% \textbf{GPT5}             & 0.9630                             & 78.5                                  & 0.6296                             & 79.8                                  & 0.8333                             & 81.7                                  & 0.9534                             & 75.9                                  & 0.8956                             & 65.1                                  & 0.8550                             & 76.2                                  \\
% \textbf{Gemini3-flash}    & 0.8889                             & 48.7                                  & 0.5556                             & 51.5                                  & 0.7963                             & 50.1                                  & 0.9265                             & 53.7                                  & 0.8342                             & 40.2                                  & 0.8003                             & 48.8                                  \\ \hline
% \end{tabular}
% \end{table*}

\begin{table*}[t]
\centering
\caption{Backbone model ablation on AgentsNet.
Token consumption is reported in thousands (K).}
\label{tab:model}

\footnotesize
\setlength{\tabcolsep}{6pt}
\renewcommand{\arraystretch}{1.10}

\resizebox{0.9\linewidth}{!}{%
\begin{tabular}{@{}l*{12}{c}@{}}
\toprule
& \multicolumn{2}{c}{\textbf{Consensus}}
& \multicolumn{2}{c}{\textbf{Leader}}
& \multicolumn{2}{c}{\textbf{Matching}}
& \multicolumn{2}{c}{\textbf{Coloring}}
& \multicolumn{2}{c}{\textbf{VertexCover}}
& \multicolumn{2}{c}{\textbf{Overall}} \\
\cmidrule(lr){2-3}
\cmidrule(lr){4-5}
\cmidrule(lr){6-7}
\cmidrule(lr){8-9}
\cmidrule(lr){10-11}
\cmidrule(lr){12-13}

\textbf{Model}
& \textbf{Score} & \textbf{Tok.}
& \textbf{Score} & \textbf{Tok.}
& \textbf{Score} & \textbf{Tok.}
& \textbf{Score} & \textbf{Tok.}
& \textbf{Score} & \textbf{Tok.}
& \textbf{Score} & \textbf{Tok.} \\
\midrule

Qwen
& 0.9630 & 57.6
& 0.5185 & 57.0
& 0.7917 & 52.4
& 0.9967 & 61.7
& 0.8342 & 43.8
& 0.8208 & 54.5 \\

DeepSeekV4-Flash
& 0.9630 & 51.0
& 0.5556 & 72.9
& 0.6435 & 64.2
& 0.9463 & 51.3
& 0.8250 & 44.6
& 0.7867 & 56.8 \\

MiniMax-2.7
& 0.9259 & 55.8
& 0.5185 & 63.4
& 0.7592 & 51.3
& 0.9177 & 55.4
& 0.7935 & 35.6
& 0.7830 & 52.3 \\

GLM-5.1
& 0.9259 & 62.9
& 0.5926 & 63.2
& 0.7314 & 56.8
& 0.9446 & 61.0
& 0.8250 & 53.1
& 0.8039 & 59.4 \\

GPT-5
& 0.9630 & 78.5
& 0.6296 & 79.8
& 0.8333 & 81.7
& 0.9534 & 75.9
& 0.8956 & 65.1
& 0.8550 & 76.2 \\

Gemini3-Flash
& 0.8889 & 48.7
& 0.5556 & 51.5
& 0.7963 & 50.1
& 0.9265 & 53.7
& 0.8342 & 40.2
& 0.8003 & 48.8 \\

\bottomrule
\end{tabular}%
}
\end{table*}

\begin{table*}[t]
\centering
\caption{Agent-count and topology ablations. Each entry reports
Score / Avg.\ Token (K).}
\label{tab:scale_topology_ablation}

\small
\renewcommand{\arraystretch}{1.15}

\begin{minipage}[t]{0.48\textwidth}
\centering
\textbf{(a) Number of Agents}

\medskip
\begin{tabular*}{\linewidth}
{@{\extracolsep{\fill}}cccc@{}}
\toprule
\textbf{\#Agents}
& \textbf{Clean}
& \textbf{Attack}
& \textbf{BRA-Audit} \\
\midrule
4
& 0.859 / 8.3
& 0.363 / 14.1
& \textbf{0.837 / 18.9} \\

8
& 0.827 / 21.3
& 0.520 / 33.4
& \textbf{0.847 / 41.9} \\

16
& 0.763 / 50.4
& 0.439 / 71.0
& \textbf{0.778 / 102.8} \\
\midrule
Overall
& 0.816 / 26.7
& 0.441 / 39.5
& \textbf{0.821 / 54.5} \\
\bottomrule
\end{tabular*}
\end{minipage}
\hfill
\begin{minipage}[t]{0.48\textwidth}
\centering
\textbf{(b) Graph Topology}

\medskip
\begin{tabular*}{\linewidth}
{@{\extracolsep{\fill}}cccc@{}}
\toprule
\textbf{Topology}
& \textbf{Clean}
& \textbf{Attack}
& \textbf{BRA-Audit} \\
\midrule
WS
& 0.803 / 26.2
& 0.451 / 38.5
& \textbf{0.846 / 55.6} \\

BA
& 0.835 / 25.2
& 0.415 / 38.1
& \textbf{0.813 / 47.8} \\

DT
& 0.811 / 28.7
& 0.456 / 42.0
& \textbf{0.804 / 60.1} \\
\midrule
Overall
& 0.816 / 26.7
& 0.441 / 39.5
& \textbf{0.821 / 54.5} \\
\bottomrule
\end{tabular*}
\end{minipage}

\end{table*}
\section{Experiment}
We evaluate BRA-Audit against recent LLM-agent auditing and multi-agent defense methods, focusing on defense effectiveness, audit cost, and robustness. Our experiments are designed to answer three research questions:

\textbf{RQ1: Defense-Cost Trade-off.} 
Can BRA-Audit achieve effective auditing while keeping the audit overhead low?

\textbf{RQ2: Robustness Across Settings.} 
How stable is BRA-Audit under different graph topologies, node importance distributions, attack locations, and backbone auditor models?

\textbf{RQ3: Budget and Component Analysis.} 
How do audit budget, audit point placement strategy, and exposure-weight components affect the overall performance of BRA-Audit?




\subsection{Experimental Setup.}

\paragraph{Dataset.} We evaluate BRA-Audit on three complementary settings. First, we use AgentsNet~\citep{grötschla2025agentsnetcoordinationcollaborativereasoning}, a multi-agent coordination benchmark with five graph-theoretic task and configurable communication topologies and node counts. This setting tests auditing under controlled multi-round agent communication. Second, we use BIG-Bench Hard (BBH)~\citep{suzgun2022challengingbigbenchtaskschainofthought} to evaluate complex logical reasoning; we randomly sample 150 questions and convert each instance into a multi-agent reasoning workflow. Third, we use the Research scenario from MultiAgentBench~\citep{zhu2025multiagentbenchevaluatingcollaborationcompetition}, an open-ended multi-agent task where agents collaborate to generate research ideas. 




\paragraph{Settings.}
We implement all multi-agent systems with AutoGen~\citep{wu2023autogen}, and instantiate agent networks with different sizes and communication topologies. In each setting, 20\% of the agents is designated as malicious and is instructed to inject incorrect claims, misleading evidence, or corrupted 
intermediate results into the collaboration process to make the task fail \textbf{(detailed attack recovery protocol is provided in technical supplement)}. We mainly use Qwen3-235B-A22B-Instruct-2507 model. We further evaluate different backbone models, graph topologies, numbers of agents, and audit budgets \(\rho\). The main metrics are task score and total E2E token consumption, which can measure defense effectiveness and cost. Each configuration is repeated three times, with the standard deviation kept below \(2\%\).



\subsection{Main Results.}

We compare BRA-Audit with five representative LLM-based methods for runtime auditing or reliability enhancement. PeerGuard performs peer-based inspection during agent interactions~\citep{fan2025peerguard}. AuditAgent is a round-level auditing baseline adapted from AgentAuditor~\citep{yang2026agentauditor}. GuardAgent converts scenario-specific guard requirements into executable checks~\citep{xiang2024guardagent}. TrinityGuard uses preassigned runtime monitors to inspect MAS execution~\citep{wang2026trinityguard}. Multi-Agent Debate improves output reliability through iterative discussion among agents~\citep{du2024multiagentdebate}. All methods use the same attack settings and backbone models.

% The attack results first confirm that LLM-MAS is vulnerable to a small number of corrupted agents. Although only one or two agents are malicious or affected by hallucination, their outputs can propagate through communication and aggregation, causing a substantial decline in the final task performance. On AgentsNet, the overall score drops from \(0.8162\) in the clean setting to \(0.4407\) under attack, a relative decrease of approximately \(46.0\%\). This result highlights the need to audit intermediate collaboration trajectories rather than relying only on the final output.

% For AgentsNet (Table~\ref{tab:agentsnet_main}), BRA-Audit recovers the overall score to \(0.8208\) with \(54.5\)K tokens. It outperforms AuditAgent, GuardAgent, and TrinityGuard, and matches the clean results on Consensus, Coloring, and VertexCover. PeerGuard and Multi-Agent Debate achieve slightly higher overall scores, but BRA-Audit reduces their token consumption by \(17.2\%\) and \(28.4\%\), respectively.

% A similar defense--cost trade-off is observed on BBH (Table~\ref{tab:BBH}). BRA-Audit improves accuracy from \(0.3933\) to \(0.8600\), while using \(19.3\%\) fewer tokens than PeerGuard and \(40.6\%\) fewer than Multi-Agent Debate. On MultiAgentBench-Research (Table~\ref{tab:BBH}), BRA-Audit reaches \(0.902\) with \(38.5\)K tokens. Its performance is close to PeerGuard and Multi-Agent Debate, while reducing token usage by \(18.5\%\) and \(36.3\%\), respectively. These results show that BRA-Audit maintains a favorable balance between defense effectiveness and runtime cost across structured and open-ended tasks.
The attack results confirm that LLM-MAS is vulnerable to a small number
of corrupted agents. Even one or two malicious or hallucinating agents
can propagate faulty outputs through communication and aggregation,
substantially degrading task performance. On AgentsNet, the overall
score falls from \(0.8162\) in the clean setting to \(0.4407\) under
attack, a relative drop of approximately \(46.0\%\). This supports auditing MAS collaboration is necessary.

On AgentsNet (Table~\ref{tab:agentsnet_main}), BRA-Audit restores the
overall score to \(0.8208\) with \(54.5\)K tokens. It outperforms
AuditAgent, GuardAgent, and TrinityGuard, and matches the clean results
on Consensus, Coloring, and VertexCover. Although PeerGuard and
Multi-Agent Debate achieve slightly higher scores, BRA-Audit uses
\(17.2\%\) and \(28.4\%\) fewer tokens, respectively.
A similar defense--cost trade-off appears on BBH and
MultiAgentBench-Research (Table~\ref{tab:BBH}). BRA-Audit improves BBH
accuracy from \(0.3933\) to \(0.8600\), using \(19.3\%\) fewer tokens
than PeerGuard and \(40.6\%\) fewer than Multi-Agent Debate. On
Research, it reaches \(0.902\) with \(38.5\)K tokens, close to both
baselines while reducing token usage by \(18.5\%\) and \(36.3\%\),
respectively. Overall, BRA-Audit balances defense effectiveness and
runtime cost across structured and open-ended tasks.



\subsection{Ablation Study}


\paragraph{Backbone Model Ablation.}
We evaluate BRA-Audit with six backbone models on AgentsNet benchmark to study its sensitivity to model capability(as shown in Table~\ref{tab:model}.) Overall scores range from \(0.7830\) to \(0.8550\), indicating stable behavior across all tested models. GPT-5 obtains the highest score of \(0.8550\), but also incurs the largest token cost at \(76.2\)K. Qwen offers a better performance--cost balance, reaching \(0.8208\) with \(54.5\)K tokens. Gemini-3 uses the fewest tokens (\(48.8\)K) while maintaining a score of \(0.8003\). GLM-5.1, DeepSeek-V4, and MiniMax-2.7 achieve \(0.8039\), \(0.7867\), and \(0.7830\), respectively. Performance differences are most visible on Leader and Matching. Overall, BRA-Audit remains effective across different backbones, while stronger models tend to improve task performance at higher token cost.



% \begin{table}[t]
% \centering
% \caption{Audit-budget ablation on BBH and MultiAgentBench-Research.
% Tokens are reported in thousands (K). Rec. denotes the absolute recovery gain over the attacked setting}
% \label{tab:budget_ablation}

% \footnotesize
% \setlength{\tabcolsep}{6pt}
% \renewcommand{\arraystretch}{1.10}

% \resizebox{\columnwidth}{!}{%
% \begin{tabular}{@{}l*{6}{c}@{}}
% \toprule
% & \multicolumn{3}{c}{\textbf{BBH}}
% & \multicolumn{3}{c}{\textbf{Research}} \\
% \cmidrule(lr){2-4}
% \cmidrule(lr){5-7}

% \textbf{\(\rho\)}
% & \textbf{Score}
% & \textbf{Rec.}
% & \textbf{Tok.}
% & \textbf{Score}
% & \textbf{Rec.}
% & \textbf{Tok.} \\
% \midrule

% \(0.2\)
% & 0.760 & 0.367 & 6.903
% & 0.757 & 0.380 & 33.2 \\

% \(0.4\), Rand.
% & 0.820 & 0.427 & 8.012
% & 0.821 & 0.444 & 37.4 \\

% \textbf{\(0.4\), BRA}
% & \textbf{0.860} & \textbf{0.467} & 8.175
% & \textbf{0.902} & \textbf{0.525} & 38.5 \\

% \(0.6\)
% & 0.873 & 0.480 & 9.568
% & 0.906 & 0.529 & 42.9 \\

% \(1.0\)
% & 0.893 & 0.500 & 11.314
% & 0.897 & 0.520 & 46.8 \\

% \bottomrule
% \end{tabular}%
% }
% \end{table}

% \begin{table}[t]
% \centering
% \caption{Audit-budget ablation with audit and rollback costs.
% All token costs are reported in thousands (K).}
% \label{tab:budget_ablation}

% \footnotesize
% \setlength{\tabcolsep}{2.6pt}
% \renewcommand{\arraystretch}{1.10}

% \resizebox{\columnwidth}{!}{%
% \begin{tabular}{@{}l*{8}{c}@{}}
% \toprule
% & \multicolumn{4}{c}{\textbf{BBH}}
% & \multicolumn{4}{c}{\textbf{Research}} \\
% \cmidrule(lr){2-5}
% \cmidrule(lr){6-9}

% \textbf{\(\rho\)}
% & \textbf{Score}
% & \textbf{Audit}
% & \textbf{Rollback}
% & \textbf{Total}
% & \textbf{Score}
% & \textbf{Audit}
% & \textbf{Rollback}
% & \textbf{Total} \\
% \midrule

% \(0.2\)
% & 0.760 & 0.737 & 2.835 & 6.903
% & 0.757 & 2.24 & 13.14 & 33.2 \\

% \(0.4\), Rand.
% & 0.820 & 2.139 & 1.757 & 8.012
% & 0.821 & 5.43 & 9.13 & 37.4 \\

% \textbf{\(0.4\), BRA}
% & \textbf{0.860} & 2.093 & 1.814 & 8.175
% & \textbf{0.902} & 5.82 & 9.38 & 38.5 \\

% \(0.6\)
% & 0.873 & 3.912 & 1.346 & 9.568
% & 0.906 & 9.11 & 7.17 & 42.9 \\

% \(1.0\)
% & 0.893 & 5.573 & 0.939 & 11.314
% & 0.897 & 11.13 & 5.64 & 46.8 \\

% \bottomrule
% \end{tabular}%
% } % closes \resizebox

% \end{table}
\begin{table}[t]
\centering
\caption{Audit-budget ablation with audit and rollback costs.
All token costs are reported in thousands (K).}
\label{tab:budget_ablation}

\footnotesize
\setlength{\tabcolsep}{2.6pt}
\renewcommand{\arraystretch}{1.10}

\resizebox{\columnwidth}{!}{%
\begin{tabular}{@{}l*{8}{c}@{}}
\toprule
& \multicolumn{4}{c}{\textbf{BBH}}
& \multicolumn{4}{c}{\textbf{Research}} \\
\cmidrule(lr){2-5}
\cmidrule(lr){6-9}

\textbf{\(\rho\)}
& \textbf{Score}
& \textbf{Audit}
& \textbf{Rollback}
& \textbf{Total}
& \textbf{Score}
& \textbf{Audit}
& \textbf{Rollback}
& \textbf{Total} \\
\midrule

\(0.2\)
& 0.760 & 0.74 & 2.84 & 6.90
& 0.757 & 2.24 & 13.14 & 33.2 \\

\(0.4\), Rand.
& 0.820 & 2.14 & 1.76 & 8.01
& 0.821 & 5.43 & 9.13 & 37.4 \\

\textbf{\(0.4\), BRA}
& \textbf{0.860} & 2.09 & 1.81 & 8.18
& \textbf{0.902} & 5.82 & 9.38 & 38.5 \\

\(0.6\)
& 0.873 & 3.91 & 1.35 & 9.57
& 0.906 & 9.11 & 7.17 & 42.9 \\

\(1.0\)
& 0.893 & 5.57 & 0.94 & 11.31
& 0.897 & 11.13 & 5.64 & 46.8 \\

\bottomrule
\end{tabular}%
}

\end{table}


\paragraph{Audit-Budget Ablation.}

We vary the audit ratio \(\rho\) on BBH and MultiAgentBench-Research to
study the defense--cost trade-off. On BBH, increasing \(\rho\) from
\(0.2\) to \(0.4\) improves the score from \(0.760\) to \(0.860\),
while total token consumption rises from \(6.9\) to \(8.18\).
Larger budgets of \(0.6\) and \(1.0\) further increase the score to
\(0.873\) and \(0.893\), but require \(9.57\) and \(11.31\)K
tokens. MultiAgentBench-Research shows a similar trend:
\(\rho=0.4\) achieves \(0.902\) with \(38.5\)K tokens, close to the
best score of \(0.906\) at \(\rho=0.6\). As \(\rho\) increases, audit
cost grows while rollback cost decreases, indicating that more frequent
auditing detects failures earlier and reduces re-execution. Under the
same budget \(\rho=0.4\), BRA-Audit outperforms random audit-point
placement by \(0.040\) on BBH and \(0.081\) on Research with comparable
token costs. Overall, \(\rho=0.4\) provides a favorable balance among
audit effectiveness, audit cost, and rollback overhead.
% We vary the audit ratio \(\rho\) on BBH and MultiAgentBench-Research to examine the defense--cost trade-off. On BBH, increasing \(\rho\) from \(0.2\) to \(0.4\) improves the task score from \(0.760\) to \(0.860\), while total token consumption increases from \(6{,}903\) to \(8{,}175\). Further increasing the budget to \(0.6\) and \(1.0\) yields scores of \(0.873\) and \(0.893\), but requires \(9{,}568\) and \(11{,}314\) tokens. A similar trend appears on MultiAgentBench-Research: \(\rho=0.4\) reaches \(0.902\) with \(38.5\)K tokens, close to the best score of \(0.906\) at \(\rho=0.6\). Moreover, under the same budget \(\rho=0.4\), BRA-Audit outperforms random audit point placement by \(0.040\) on BBH and \(0.081\) on Research with comparable token costs. Overall, \(\rho=0.4\) provides a favorable balance between audit effectiveness and runtime overhead.


% \paragraph{Node-Count Ablation.}
% We evaluate BRA-Audit on AgentsNet graphs with 4, 8, and 16 agents, as shown in Table~\ref{tab:scale_topology_ablation}a. BRA-Audit obtains scores of \(0.8370\), \(0.8473\), and \(0.7781\), respectively, and consistently recovers a large portion of the performance lost under attack. The 8-agent setting gives the highest score, while the 16-agent setting is more challenging because more agents create longer communication paths and larger audit contexts. Token consumption increases from \(18.9\)K to \(41.9\)K and \(102.8\)K as the network grows. Compared with the attacked setting, BRA-Audit improves the score by \(0.4737\), \(0.3274\), and \(0.3391\) across the three sizes. These results show that BRA-Audit remains effective as the network scales, although larger systems require substantially more auditing resources.

% \paragraph{Topology Ablation.}
% We test BRA-Audit on the three graph families used by AgentsNet: Watts--Strogatz small-world graphs, Barabási--Albert preferential-attachment graphs, and Delaunay triangulations~\citep{grötschla2025agentsnetcoordinationcollaborativereasoning}. AS shown in Table~\ref{tab:scale_topology_ablation}b, these topologies respectively represent highly clustered networks with short paths, hub-dominated scale-free networks, and spatially local communication structures. BRA-Audit achieves scores of \(0.8458\), \(0.8127\), and \(0.8040\), compared with attacked scores of \(0.4511\), \(0.4153\), and \(0.4558\). The small-world topology gives the highest task score, whereas the preferential-attachment topology uses the fewest tokens at \(47.8\)K. The Delaunay setting requires \(60.1\)K tokens and records a score of \(0.8040\). Overall, BRA-Audit provides substantial recovery across different communication structures, although topology affects both task performance and audit overhead.

\paragraph{Node-Count Ablation.}
We evaluate BRA-Audit on AgentsNet graphs with 4, 8, and 16 agents
(Table~\ref{tab:scale_topology_ablation}a). It achieves scores of
\(0.8370\), \(0.8473\), and \(0.7781\), improving over the attacked
setting by \(0.4737\), \(0.3274\), and \(0.3391\), respectively.
Token consumption rises from \(18.9\)K to \(41.9\)K and \(102.8\)K as
larger graphs introduce longer communication paths and audit contexts.
BRA-Audit therefore remains effective as the system scales, although
larger networks require substantially more auditing resources.

\paragraph{Topology Ablation.}
We evaluate Watts--Strogatz, Barabási--Albert, and Delaunay graphs,
which represent clustered small-world, hub-dominated, and spatially
local communication structures, respectively
\citep{grötschla2025agentsnetcoordinationcollaborativereasoning}.
As shown in Table~\ref{tab:scale_topology_ablation}b, BRA-Audit obtains
scores of \(0.8458\), \(0.8127\), and \(0.8040\), compared with attacked
scores of \(0.4511\), \(0.4153\), and \(0.4558\). The small-world graph
achieves the best score, while the preferential-attachment graph uses
the fewest tokens. Overall, BRA-Audit remains effective across diverse
communication topologies.

\begin{table}[t]
\centering
\caption{False-positive results under different audit ratios.}
\label{tab:false_positive}
\footnotesize
\renewcommand{\arraystretch}{1.12}

\begin{tabular*}{\columnwidth}{
@{\extracolsep{\fill}}
c
cc
cc
@{}
}
\toprule
& \multicolumn{2}{c}{\textbf{BBH}}
& \multicolumn{2}{c}{\textbf{Research}} \\
\cmidrule(lr){2-3}
\cmidrule(lr){4-5}
\(\boldsymbol{\rho}\)
& \textbf{FP}
& \textbf{FPR}
& \textbf{FP}
& \textbf{FPR} \\
\midrule
\(0.2\) & 0  & \(0.00\%\)  & 0  & \(0.00\%\)  \\
\(0.4\) & 7  & \(4.67\%\)  & 6  & \(6.00\%\)  \\
\(1.0\) & 18 & \(12.00\%\) & 16 & \(16.00\%\) \\
\bottomrule
\end{tabular*}
\end{table}

\paragraph{False-Positive Analysis.}
We assess false positives on clean BBH and MultiAgentBench-Research at varying audit ratios (Table~\ref{tab:false_positive}). No FP occur at \(\rho=0.2\). At \(\rho=0.4\), FP rates are
\(4.67\%/6.00\%\) for BBH/Research, rising to
\(12.00\%/16.00\%\) under full auditing. Thus, excessive auditing can introduce unnecessary errors on benign trajectories. Most executions contain neither attacks nor severe hallucinations, so auditing every interaction incurs substantial redundant cost. Alongside task performance, BRA-Audit at \(\rho=0.4\) empirically
lies on a favorable Pareto frontier balancing defense effectiveness, FP risk, and token consumption.
% We evaluate false positives on clean BBH and MultiAgentBench-Research instances under different audit ratios (Table~\ref{tab:false_positive}). No false positives occur at \(\rho=0.2\). At \(\rho=0.4\), the false-positive rates are \(4.67\%\) on BBH and \(6.00\%\) on Research, increasing to \(12.00\%\) and \(16.00\%\) under full auditing. These results show that excessive auditing can introduce unnecessary errors even on benign trajectories. Since most executions do not contain attacks or severe hallucinations, auditing every interaction also incurs substantial redundant cost. Combined with the task-performance results, BRA-Audit at \(\rho=0.4\) empirically lies on a favorable Pareto frontier between defense effectiveness, false-positive risk, and token consumption.

% \paragraph{Weight-Component Ablation.} we also do the weight component ablations, results show that both \(\mathrm{AuditGap}\) and \(\mathrm{TopoImpact}\) improve BRA-Audit. Removing
% \(\mathrm{AuditGap}\) lowers the overall score from \(0.8208\) to \(0.7578\), while removing \(\mathrm{TopoImpact}\) reduces it to \(0.7932\); both variants also consume more tokens. These results
% confirm their complementary roles in temporal coverage and propagation-aware scheduling. \textbf{Due to space constraints, detailed results and analysis are provided in technical supplement.}

\paragraph{Weight-Component Ablation.} We further ablate the two components of the exposure weight. Removing \(\mathrm{AuditAge}\) decreases the overall score from \(0.821\) to \(0.758\), while removing \(\mathrm{ClosedImpact}\) reduces it to \(0.793\). Both variants also increase total token consumption, showing that verification staleness and closed downstream influence jointly improve audit-point selection. Detailed results and analysis are provided in the \textbf{technical supplement}.



\section{Conclusion}

% In this paper, we presented BRA-Audit to address the trade-off between defense effectiveness and runtime cost in LLM multi-agent systems. BRA-Audit models MAS communication as a dependency graph and formulates runtime auditing as a cumulative-exposure audit point placement problem. By selecting a limited number of audit points, it shortens audit trajectories, limits the propagation of unchecked exposure, and keeps rollback localized. Its scheduling objective combines audit gaps with topological impact, directing audit resources toward long-unaudited and influential regions. Experiments on structured coordination, complex reasoning, and open-ended multi-agent tasks show that BRA-Audit provides a favorable balance between task recovery and token cost. Results across different models, graph structures, network sizes, and audit budgets further demonstrate its practical stability. We believe BRA-Audit can contribute to the development of more trustworthy and cost-aware MAS.

In this paper, we presented BRA-Audit, a cost-aware runtime defense for LLM MAS. It models MAS communication as a dependency graph and formulates auditing as cumulative-exposure audit point placement. By combining audit gaps with topological impact, BRA-Audit selects limited audit points in long-unaudited and influential regions, shortens audit contexts, limits unchecked propagation, and localizes rollback. Experimental results on structured coordination, complex reasoning, and open-ended tasks support a favorable balance between task recovery and token cost. demonstrates a favorable Pareto-optimal trade-off between task performance and auditing cost.  We believe BRA-Audit can contribute to the development of more trustworthy and cost-aware MAS.

\bibliography{aaai2027}

% Check whether the conference requires a reproducibility checklist to be included in the paper.
% If so, you can uncomment the following line and ajust the path to include it.
% \input{ReproducibilityChecklist.tex}

\end{document}
